% year: תשע"ז
% type: מבחן
% semester: א
% moed: א
% number: 4א
% points: 12
% categories: func-analysis
% source: exams/מבחנים/תשעז/מועד א תשעז.pdf

\begin{question}
חקרו את הפונקציה $y=\ln(e^x+e^{-2x})$ וסרטטו את הגרף שלה.
\underline{צריך למצוא}: תחום הגדרה, אסימפטוטות, תחומי עליה וירידה, נקודות קיצון, תחומי קמירות וקעירות.
\end{question}

\begin{hint}
לאסימפטוטות: כתבו $\ln(e^x+e^{-2x})=x+\ln(1+e^{-3x})=-2x+\ln(1+e^{3x})$.
\end{hint}

\begin{solution}
\textbf{תחום הגדרה:} $e^x+e^{-2x}>0$ לכל $x$, ולכן הפונקציה מוגדרת (ורציפה וגזירה) על כל $\R$.

\textbf{אסימפטוטות:} אין אסימפטוטות אנכיות, כי הפונקציה רציפה על כל $\R$.
עבור $x\to+\infty$:
\[
y=\ln\big(e^x(1+e^{-3x})\big)=x+\ln(1+e^{-3x}),\qquad y-x=\ln(1+e^{-3x})\xrightarrow[x\to+\infty]{}0,
\]
ולכן $y=x$ אסימפטוטה משופעת ב-$+\infty$.
עבור $x\to-\infty$:
\[
y=\ln\big(e^{-2x}(1+e^{3x})\big)=-2x+\ln(1+e^{3x}),\qquad y+2x\xrightarrow[x\to-\infty]{}0,
\]
ולכן $y=-2x$ אסימפטוטה משופעת ב-$-\infty$. אין אסימפטוטות אופקיות. בנוסף $y-x>0$ ו-$y+2x>0$, כלומר הגרף נמצא מעל שתי האסימפטוטות.

\textbf{עליה וירידה:}
\[
y'=\frac{e^x-2e^{-2x}}{e^x+e^{-2x}}.
\]
המכנה חיובי, ו-$e^x-2e^{-2x}>0\iff e^{3x}>2\iff x>\frac{\ln2}{3}$.
לכן $y$ יורדת ב-$\left(-\infty,\frac{\ln 2}{3}\right)$ ועולה ב-$\left(\frac{\ln2}{3},\infty\right)$.

\textbf{קיצון:} ב-$x_0=\frac{\ln2}{3}$ יש מינימום (מוחלט). שם $e^{x_0}=2^{1/3}$, $e^{-2x_0}=2^{-2/3}$, ולכן
\[
y(x_0)=\ln\left(2^{1/3}+2^{-2/3}\right)=\ln\left(3\cdot2^{-2/3}\right)=\ln3-\frac23\ln2\approx0.64.
\]

\textbf{קמירות:} לפי כלל המנה,
\[
y''=\frac{(e^x+4e^{-2x})(e^x+e^{-2x})-(e^x-2e^{-2x})^2}{(e^x+e^{-2x})^2}
=\frac{9e^{-x}}{(e^x+e^{-2x})^2}>0,
\]
(במונה: $e^{2x}+5e^{-x}+4e^{-4x}-e^{2x}+4e^{-x}-4e^{-4x}=9e^{-x}$). לכן הפונקציה קמורה (מחייכת, $\cup$) על כל $\R$, ואין נקודות פיתול.

\textbf{סרטוט:} הגרף יורד מלמעלה לאורך האסימפטוטה $y=-2x$ (מעליה), מגיע למינימום בנקודה $\left(\frac{\ln2}{3},\ \ln3-\frac23\ln2\right)\approx(0.23,\,0.64)$, ועולה לאורך האסימפטוטה $y=x$ (מעליה); כולו קמור. חיתוך עם ציר $y$: $y(0)=\ln2$. אין חיתוך עם ציר $x$ (כי $y\ge y(x_0)>0$).
\end{solution}
